%%%%%%%%%%%%%%%%%%%%%%%%%%%%%%%%%%%%%%%%%%%%%%%%%%%%%%%
% % Lines starting with % are comments, which are ignored.
% % This is a handy way of indicating the date and version of
% % your document, to wit:
% %
% % LaTeX sample file
% % Modified March, 2002
% %
%%%%%%%%%%%%%%%%%%%%%%%%%%%%%%%%%%%%%%%%%%%%%%%%%%%%%%%
\DocumentMetadata{
  % PDF standards: ONE key whose value is a comma-separated list --
  % UA-2 (accessibility) plus A-4f (archival; the archival companion of
  % UA-2). Do NOT use separate pdfstandard=... keys: with ordinary
  % key-value parsing only the last one takes effect.
  pdfstandard={UA-2,A-4f},
  lang=en-US,
  tagging=on,
  tagging-setup={
    math/setup=mathml-SE
  }
}

\documentclass{ltx-talk}
% lecture-macros.tex -- shared package setup for lecture-template.tex
% (ltx-talk slides / handouts).
%
% NOTE ON PATHS: every template in this repository is stored one level up
% from where it is compiled (see README.md). All relative \input paths
% below are resolved against the folder the document is compiled in,
% which sits one directory above this file.

% --- Tables, bookmarks, strikethrough ------------------------------------
% booktabs: professional table rules (\toprule, \midrule, \bottomrule)
% bookmark: PDF outline entries (PDF bookmarks)
% ulem[normalem]: \sout strikethrough that survives math mode
\usepackage{booktabs}
\usepackage{bookmark}
\usepackage[normalem]{ulem}
\usepackage{multicol}

% --- Math ------------------------------------------------------------
% amsmath: equation environments and numbering
% mathrsfs: \mathfrak  |  stmaryrd: extra symbols
\usepackage{amsmath}
\usepackage{mathrsfs}
\usepackage{stmaryrd}

% --- Layout -------------------------------------------------------------
% float: [H] float placement
% fontspec: Unicode/OpenType fonts (compile with XeLaTeX or LuaLaTeX)
\usepackage{float}
\usepackage{fontspec}

% --- Colors ---------------------------------------------------------------
% dvipsnames: extended named color palette (the Accessible* colors below
% are built on it);
% table: \rowcolors / \cellcolor for shading tables.
% NOTE: xcolor may only be loaded ONCE per document -- a second
% \usepackage{xcolor} with different options is a fatal "Option clash"
% error, so ALL of its options are listed here.
\usepackage[dvipsnames,table]{xcolor}

% --- Multimedia -----------------------------------------------------------
% animate: embeds audio/video (and legacy Flash) media in a frame.
\usepackage{animate}

% --- Bibliography ----------------------------------------------------------
% biblatex (biber backend): citations for lecture notes.
\usepackage[
  backend=biber,
  style=alphabetic,
  url=true,
  doi=true,
  isbn=true,
  backref=true
]{biblatex}

% --- Colored boxes ----------------------------------------------------------
% talkthemetcolorbox: the colored-box theme for ltx-talk slides
% (shadow/rounded/shaded blocks, showtitle, titlepage styles).
\usepackage[
  shadow,
  rounded,
  shaded,
  blocks,
  showtitle,
  titlepage
]{talkthemetcolorbox}

% --- "theindex" helper environment ------------------------------------------
% Provides a lightweight theindex environment (an itemize) and the
% \@idxitem/\subitem/\subsubitem/\indexspace macros that LaTeX's standard
% theindex machinery expects, so hand-written index pages indent correctly.
\makeatletter
\newenvironment{theindex}
  {\begin{itemize}}
  {\end{itemize}}
\providecommand\@idxitem{\par\hangindent 40\p@}
\providecommand\subitem{\@idxitem \hspace*{20\p@}}
\providecommand\subsubitem{\@idxitem \hspace*{30\p@}}
\providecommand\indexspace{\par \vskip 10\p@ \@plus5\p@ \@minus3\p@\relax}
\makeatother

% --- Index -----------------------------------------------------------------
% imakeidx: builds the printed index.
% The redefinition below appends "|hyperpage" to every index entry so that,
% in the PDF, each index entry is a link to its page.
\usepackage{imakeidx}
\makeatletter
\let\imki@orig@wrindexentrysplit\imki@wrindexentrysplit
\renewcommand\imki@wrindexentrysplit[3]{%
  \imki@orig@wrindexentrysplit{#1}{#2|hyperpage}{#3}%
}

\makeatletter
% --- Section titles: capture + one bookmark per logical section -------------
% \section is redefined so that:
%   1. the section title is captured into the plain-TeX register
%      \mytalksectitletoks (reprinted verbatim on the auto divider frame,
%      see the cmd/section/after hook below -- plain-TeX storage is
%      catcode-safe, unlike expanding the title in a hook), and
%   2. exactly ONE depth-1 PDF bookmark is created per \section call
%      (hyperref's own bookmark is suppressed via \Hy@writebookmark).
\DeclareCommandCopy{\OriginalSectionForBookmarks}{\section}
\newtoks\mytalksectitletoks

\RenewDocumentCommand{\section}{s o m}{%
  \begingroup
    \def\Hy@writebookmark##1##2##3##4##5{}%
    \global\mytalksectitletoks={#3}%
    \IfBooleanTF{#1}
      {\phantomsection\IfNoValueTF{#2}{\OriginalSectionForBookmarks*{#3}}{\OriginalSectionForBookmarks*[#2]{#3}}}%
      {\IfNoValueTF{#2}{\OriginalSectionForBookmarks{#3}}{\OriginalSectionForBookmarks[#2]{#3}}}%
  \endgroup
  \IfBooleanTF{#1}
    {\IfNoValueTF{#2}{\pdfbookmark[1]{#3}{section-star.\thepage.\number\inputlineno}}{\pdfbookmark[1]{#2}{section-star.\thepage.\number\inputlineno}}}%
    {\IfNoValueTF{#2}{\pdfbookmark[1]{Section \thesection: #3}{section.\thesection}}{\pdfbookmark[1]{Section \thesection: #2}{section.\thesection}}}%
}
\makeatother

% --- One PDF bookmark per logical frame (NOT per overlay) ---------------------
% ltx-talk writes one PDF page per overlay, so a naive frametitle hook would
% register a new bookmark for every overlay of the same frame. The
% \g_slideutils_frame_bookmark_int register remembers the last frame that
% got a bookmark, so subsequent overlays of the same frame are skipped.
% This is also what keeps BeamerPresenter's cumulative annotations working:
% one logical frame must stay one logical unit in the PDF outline.
\ExplSyntaxOn
\int_new:N \g_slideutils_frame_bookmark_int

\AddToHook{cmd/frametitle/after}{%
  \tl_if_empty:NF \g__talk_frame_title_tl
    {
      \ifnum\value{frame}=\g_slideutils_frame_bookmark_int\relax
      \else
        \int_gset:Nn \g_slideutils_frame_bookmark_int { \value{frame} }%
        \pdfbookmark[2]{\tl_use:N \g__talk_frame_title_tl}
          {frame.\int_use:N \g_slideutils_frame_bookmark_int}
      \fi
    }
}
\ExplSyntaxOff
\makeatother

% --- Printed index + theme for the colored boxes -------------------------------
% \makeindex[...]: two-column index titled "Index", listed in the TOC.
% The two EditInstance calls style the standard header/footer instances
% of talkthemetcolorbox (cyan background; footer order title/author/
% date/framenumber).
\makeindex[columns=2, title=Index, intoc]\EditInstance{header}{std}{
  background-color = cyan,
  color            = black,
  height           = 0.75cm
}
\EditInstance{footer}{std}{
  background-color = cyan,
  element-order = {title, author, date, framenumber},
}

% --- Auto section divider frame ---------------------------------------------------
% After each \section, insert a centered, large, structure-colored frame
% that prints the section title (from \mytalksectitletoks).
\AddToHook{cmd/section/after}{%
  \begin{frame}
    \centering
    \color{structure}
    \LARGE
    \the\mytalksectitletoks
  \end{frame}%
}

% --- Tagging sockets (tagged PDF structure) ---------------------------------------
% ltx-talk exposes tagging "sockets" where the PDF structure can attach
% section titles; register one and leave it unassigned (none) so the
% class's own tagging keeps control.
\ExplSyntaxOn
\AtBeginDocument{
  \NewTaggingSocketPlug { talk / sec / title } { none } { }
  \AssignTaggingSocketPlug { talk / sec / title } { none }
}
\ExplSyntaxOff

% --- Accessibility roles ----------------------------------------------------------------
% Map \frametitle to the H1 heading role in the tagged PDF, so screen
% readers treat slide titles as headings.
\tagpdfsetup{role / new-tag = frametitle / H1}

\makeatletter
% --- Page label = frame number -------------------------------------------------------
% Redefine what feeds the PDF page label to track the frame counter,
% not the physical page counter. Replace \theframenumber below with
% whatever internal macro/property ltx-talk actually uses to produce
% the "framenumber" footer element — confirm via ltx-talk-frame.dtx.
\renewcommand\thepage{\arabic{frame}}
\makeatother

\addbibresource{refs.bib}

% --- "Problem Time!" slides ------------------------------------------------------------------
% \ProblemSlide{problem}{solution} emits two frames: the problem, then
% its solution after a \pause, each numbered by the problem counter.
\newcounter{problemcounter}
\setcounter{problemcounter}{0}
\newcommand{\ProblemSlide}[2]{
  \stepcounter{problemcounter}
  \begin{frame}
    \frametitle{Problem Time!}
    Problem \theproblemcounter: #1
  \end{frame}

  \begin{frame}
    \frametitle{Problem Time!}
    Solution to Problem \theproblemcounter: \pause

    #2
  \end{frame}
}
% \highlight[<colour>]{<stuff>}
% from https://tex.stackexchange.com/questions/33401/a-version-of-colorbox-that-works-inside-math-environments
\newcommand{\highlight}[2][yellow]{\mathchoice%
  {\colorbox{#1}{$\displaystyle#2$}}%
  {\colorbox{#1}{$\textstyle#2$}}%
  {\colorbox{#1}{$\scriptstyle#2$}}%
  {\colorbox{#1}{$\scriptscriptstyle#2$}}}%

% --- Fonts ------------------------------------------------------------------------------
% Lato (text) + LeteSansMath (math): open, accessible typefaces.
\setmainfont{Lato}[
  Ligatures=TeX,
  Numbers={Lining,Monospaced}
]
\setsansfont{Lato}[
  Ligatures=TeX,
  Numbers={Lining,Monospaced}
]
\setmathfont{LeteSansMath.otf}\usepackage{emoji}

% --- Accessible colors ---------------------------------------------------------------------
% WCAG-contrast-safe variants of red/green for text on white:
% red!80!black and green!40!black both clear the 4.5:1 requirement.
% AltYellow: soft background for highlight boxes.
\colorlet{AccessibleRed}{red!80!black}
\colorlet{AccessibleGreen}{green!40!black}
\definecolor{AltYellow}{RGB}{255,245,170}

% --- Links ---------------------------------------------------------------------------------
% hyperref + link colors: every link uses AccessibleGreen (7.2:1 on
% white), so internal, citation, and URL links all pass WCAG 1.4.3.
\usepackage{hyperref}
\hypersetup{
  colorlinks=true,
  linkcolor=AccessibleGreen,
  urlcolor=AccessibleGreen,
  citecolor=AccessibleGreen,
}

% --- Shared course macros ---------------------------------------------------------------------
% Root macros.tex: math shorthands used across all course documents
% (\RR, \NN, \ZZ, \QQ, \grad, \curl, \tr, \der, \pder, ...).
\newcommand{\RR}{\mathbb{R}}
\newcommand{\NN}{\mathbb{N}}
\newcommand{\ZZ}{\mathbb{Z}}
\newcommand{\QQ}{\mathbb{Q}}
\newcommand{\vv}{\vec{v}}
\newcommand{\ww}{\vec{w}}
\newcommand{\xx}{\vec{x}}
\newcommand{\yy}{\vec{y}}
\newcommand{\zz}{\vec{z}}
\newcommand{\uu}{\vec{u}}
\newcommand{\basis}[1]{\vec{e_{#1}}}
\newcommand{\id}{\mathrm{id}}
\newcommand{\CC}{\mathbb{C}}
\renewcommand{\div}{\operatorname{div}}
\newcommand{\grad}{\operatorname{grad}}
\newcommand{\curl}{\operatorname{curl}}
\newcommand{\tr}{\operatorname{tr}}
\newcommand{\der}[2]{\dfrac{d #1}{d #2}}
\newcommand{\pder}[3][]{\dfrac{\partial^{#1} #2}{\partial #3}}


% % Title and author(s)
%%%%%%%%%%%%%%%%%%%%%%%%%%%%%%%%%%%%%%%%%%%%%%%%%%%%%%%
\title{Lecture 10 - Product/Quotient Rule, Trig Derivatives}
\author{Garrett Credi}
\date{10/13/2026}
%%%%%%%%%%%%%%%%%%%%%%%%%%%%%%%%%%%%%%%%%%%%%%%%%%%%%%%
% %
% % The next command allows your in import encapsulated
% % postscript files, .epsf or .eps files, which
% % contain vector graphic image data.
% %
%%%%%%%%%%%%%%%%%%%%%%%%%%%%%%%%%%%%%%%%%%%%%%%%%%%%%%%
% %   The first thanks indicates your affiliation
% %
% %  Just the name here.
% %
% % Your mailing address goes at the end.
% %
% % \thanks is also how you indicate grant support
% %
%%%%%%%%%%%%%%%%%%%%%%%%%%%%%%%%%%%%%%%%%%%%%%%%%%%%%%%

%%%%%%%%%%%%%%%%%%%%%%%%%%%%%%%%%%%%%%%%%%%%%%%%%%%%%%%
% $ The next command ensures the margins are 1 inch on
% % all sides.
% %
%%%%%%%%%%%%%%%%%%%%%%%%%%%%%%%%%%%%%%%%%%%%%%%%%%%%%%%
\begin{document}
\begin{frame}
  \maketitle
\end{frame}

\begin{frame}
  \frametitle{Exam 1 Recap}

  \begin{enumerate}
    \item All grades for Exam 1 released, regrade requests open until next Monday.
    \item After Monday, curve will be determined
    \item Recap: Essential Discontinuity, Vertical Tangent Singularity, Vertical Asymptotes
  \end{enumerate}

  Also, \href{https://forms.gle/GeHKGnT2nfTgQTGw8}{Anonymous Feedback Form}.
\end{frame}

\begin{frame}
  \frametitle{Quiz Time!}

  Let's get into neat rows, spread out, and slay this quiz :)
\end{frame}

\section{Product and Quotient Rule}

\begin{frame}
  \frametitle{Long Term Goals}

  Recall we're trying to build a \textit{synthetic} approach to studying derivatives: \pause

    Atomic Rules: \pause
  \begin{multicols}{2}
    \begin{itemize}
      \item Power Rule: $\dfrac{d}{dx} x^n$ 
      \item Trigonometric Rules: $\dfrac{d}{dx} \sin(x)$ etc.
      \item Exponential/Logarithm Rules: $\dfrac{d}{dx} e^x$, $\dfrac{d}{dx} \ln(x)$
    \end{itemize}
  \end{multicols}

  Combination Rules: \pause
  \begin{multicols}{2}
    \begin{itemize}
      \item Sum/Difference Rule: $\dfrac{d}{dx} (f(x) \pm g(x))$ 
      \item Product Rule: $\dfrac{d}{dx} f(x)g(x)$ 
      \item Chain Rule: $\dfrac{d}{dx} g(f(x))$ 
      \item Constant Multiple Rule: $\dfrac{d}{dx} c \cdot f(x)$
      \item Quotient Rule: $\dfrac{d}{dx} \dfrac{f(x)}{g(x)}$ 
      \item Inverse Function Rule: $\dfrac{d}{dx} f^{-1}(x)$
    \end{itemize}
  \end{multicols}
\end{frame}

\begin{frame}
  \frametitle{Product Rule I}

  We want to determine what $\dfrac{d}{dx} (f(x) \cdot g(x))$
  is, hopefully in terms of simpler derivatives involving $f(x), g(x)$. \pause

  Is it just $f'(x)g'(x)$? \pause \textcolor{red}{NO NO NO NO NO!}

  \begin{figure}[H]
    \centering
    \includegraphics[height=0.6\textheight]{inde-no.png}
    \caption{My one wish willow is for you to not compute derivates of products like this!}
    \label{fig:inde-no-gif}
  \end{figure}
\end{frame}

\begin{frame}
  \frametitle{Product Rule II}

  Well, then what is it? \pause

  Approach 1: Linearization! \pause
  We motivated the other rules by examining what the operation did to linear functions first - let's try that here! \pause

  Linearization of $f(x)$ at $x = x_0$: $f(x_0) + f'(x_0)(x - x_0)$,

  Linearization of $g(x)$ at $x = x_0$: $g(x_0) + g'(x_0)(x - x_0)$, \pause

  Linearization of $f(x)g(x)$ at $x = x_0$:
  \begin{align}
    f(x)g(x) &\approx (f(x_0) + f'(x_0)(x - x_0)) \cdot (g(x_0) + g'(x_0)(x - x_0)) \\ \pause
             &= f(x_0)g(x_0) + f'(x_0)g(x_0)(x - x_0) + f(x_0)g'(x_0)(x - x_0) + f'(x_0)g'(x_0)(x - x_0)^2 \\ \pause
             &= f(x_0)g(x_0) + (f'(x_0)g(x_0) + f(x_0)g'(x_0))(x - x_0) + f'(x_0)g'(x_0)\color{red}{(x-x_0)^2}
  \end{align} \pause

  So why is it \textcolor{red}{NOT} $f'(x)g'(x)$? \pause Because that's the coefficient of something involving a \textit{quadratic} term, not linear! \pause

  The derivative would be the slope on the \textit{linear} term.
\end{frame}

\begin{frame}
  \frametitle{Product Rule III}

  \vspace{\stretch{0.2}}
  \begin{alertblock}{The Product Rule}\index{derivative!product rule}
    \begin{equation}
      \dfrac{d}{dx}(f(x) \cdot g(x)) = f'(x)g(x) + f(x)g'(x)
    \end{equation}
  \end{alertblock}

  \vspace{\stretch{1}}
\end{frame}

\begin{frame}
  \frametitle{Product Rule Examples I}

  Example 1: What is the derivative of:
  \begin{equation}
    x^2 e^x + \dfrac{2^x}{x}?
  \end{equation} \pause

  We compute: \pause
  \begin{align}
    \dfrac{d}{dx}(x^2 e^x + \dfrac{2^x}{x}) &= \pause \dfrac{d}{dx}(x^2 e^x) + \dfrac{d}{dx}(\dfrac{2^x}{x}) \\ \pause
                                            &= (\dfrac{d}{dx} x^2) \cdot (e^x) + (x^2) \cdot (\dfrac{d}{dx} e^x) + \dfrac{d}{dx}(2^x \cdot \dfrac{1}{x}) \\ \pause
                                            &= (2x) \cdot e^x + x^2 \cdot e^x + (\dfrac{d}{dx} 2^x) \cdot (\dfrac{1}{x}) + (2^x) \cdot (\dfrac{d}{dx} \dfrac{1}{x}) \\ \pause
                                            &= 2xe^x + x^2 e^x + \ln(2) 2^x \cdot \dfrac{1}{x} + 2^x \cdot \dfrac{-1}{x^2}
  \end{align}
\end{frame}

\begin{frame}
  \frametitle{Product Rule Examples II}
  For a more applied example: \pause
  Let's suppose you drive $m(t)$ miles per day,
  in a car whose fuel efficiency is given by $e(t)$ gallons/mile
  where the price of fuel is given by $p(t)$ dollars/gallon,
  $t$ is measured in days. \pause

  Suppose we've measured the relevant quantities at:
  \begin{gather}
    m(1) = 50 mi, \quad e(1) = \dfrac{1}{25} \dfrac{gal}{mi}, \quad p(1) = 4.15 \dfrac{\$}{gal} \\
    m'(1) = 0 \dfrac{mi}{day}, \quad e'(1) = \dfrac{1}{1000} \dfrac{\frac{gal}{mi}}{day}, \quad p'(1) = 0.01 \dfrac{\frac{\$}{gal}}{day}
  \end{gather} \pause

  What is the rate of change of the cost to drive your car on day 1? \pause
\end{frame}
\begin{frame}
  \frametitle{Product Rule Examples III}
  Step 1: Cost Function. \pause $C(t) = m(t) \cdot e(t) \cdot p(t)$. \pause

  Step 2: Derivative. \pause
  \begin{align}
    C'(t) &= \dfrac{d}{dt}(m(t) \cdot e(t) \cdot p(t)) = \dfrac{d}{dt} ((m(t)e(t)) \cdot p(t)) \\ \pause
          &= \dfrac{d}{dt}(m(t) \cdot e(t)) \cdot p(t) + (m(t) \cdot e(t)) \cdot p'(t) \\ \pause
          &= m'(t) e(t) p(t) + m(t) e'(t) p(t) + m(t) e(t) p'(t)
  \end{align} \pause
\end{frame}
\begin{frame}
  \frametitle{Product Rule Examples IV}
  Step 3: Compute. \pause
  \begin{align}
    C'(1) &= m'(1) e(1) p(1) + m(1) e'(1) p(1) + m(1) e(1) p'(1) \\ \pause
          &= 0 \cdot \dfrac{1}{25} \dfrac{gal}{mi} \cdot 4.15 \dfrac{\$}{gal} \\
          &\quad+ 50 mi \cdot \dfrac{1}{1000} \dfrac{\frac{gal}{mi}}{day} \cdot 4.15 \dfrac{\$}{gal} \\
          &\quad+ 50 mi \cdot \dfrac{1}{25} \dfrac{gal}{mi} \cdot 0.01 \dfrac{\frac{\$}{gal}}{day} \\ \pause
          &= 0 \dfrac{\$}{day} + 0.2075 \dfrac{\$}{day} + 0.02 \dfrac{\$}{day} \\ \pause
          &= 0.2275 \dfrac{\$}{day}
  \end{align}
\end{frame}

\begin{frame}
  \frametitle{Quotient Rule I}

  Similarly, we'd want to determine $\dfrac{d}{dx} \dfrac{f(x)}{g(x)}$,
  again in terms of simpler derivatives involving $f(x)$, $g(x)$. \pause

  Is it just $\dfrac{f'(x)}{g'(x)}$? \pause \textcolor{red}{SUPER NO! NO! NO!} 

  \begin{figure}[H]
    \centering
    \includegraphics[height=0.4\textheight]{inde-no.png}
    \caption{I'm amending my one wish willow to also include you not computing derivatives of quotients like this.}
    \label{fig:inde-no-png}
  \end{figure}
\end{frame}

\begin{frame}
  \frametitle{Quotient Rule II}

  We can determine what it actually is by realizing that:
  \begin{equation}
    \dfrac{f(x)}{g(x)} = f(x) \cdot \dfrac{1}{g(x)}
  \end{equation} \pause
  and since we know how to linearize $f(x)$ and how to linearize a product,
  \pause
  we need only understand how to linearize something like $\dfrac{1}{g(x)}$. \pause

  First, we note that for the derivative of $\dfrac{f(x)}{g(x)}$ at $x = x_0$ to exist, \pause the function $\dfrac{f(x)}{g(x)}$ must be continuous at $x = x_0$, \pause so $g(x_0) \neq 0$.

  This allows us to approximate $\dfrac{1}{g(x)} \approx \dfrac{1}{g(x_0) + g'(x_0)(x - x_0)}$.
\end{frame}

\begin{frame}
  \frametitle{Quotient Rule III}
  Since
  \begin{equation}
    \dfrac{1}{g(x)} \approx \dfrac{1}{g(x_0) + g'(x_0)(x - x_0)} \pause = \dfrac{1}{(g(x_0) - g'(x_0)x_0) + g'(x_0)x}
  \end{equation} \pause

  we can focus on understanding how to linearize functions like
  \begin{equation}
    \dfrac{1}{a + bx}.
  \end{equation}

  We compute: \pause
  \begin{align}
    \dfrac{d}{dx} \dfrac{1}{a + bx} &= \dfrac{d}{dx} \dfrac{1}{b} \dfrac{1}{x + \dfrac{a}{b}}  \pause
                                    = \dfrac{1}{b} \dfrac{d}{dx} \dfrac{1}{x + \dfrac{a}{b}} \pause \\
                                    &= \dfrac{1}{b} \dfrac{-1}{(x + \dfrac{a}{b})^2}  \pause
                                    = \dfrac{1}{b} \dfrac{-1}{x^2 + \frac{2a}{b}x + \frac{a^2}{b^2}} \pause \\
                                    &= \dfrac{-1}{bx^2 + 2a x + \frac{a^2}{b}}  \pause
                                    = \dfrac{-b}{b^2x^2 + 2abx + a^2}  \pause
                                    = \dfrac{-b}{(a + bx)^2}
  \end{align} \pause
\end{frame}

\begin{frame}
  \frametitle{Quotient Rule IV}
  And so,

  \begin{equation}
    \dfrac{1}{g(x)} \approx \dfrac{1}{g(x_0)} + \dfrac{-g'(x_0)}{g(x_0)^2} (x - x_0)
  \end{equation} \pause

  meaning that: \pause

  \begin{align}
    \dfrac{f(x)}{g(x)} &= f(x) \cdot \dfrac{1}{g(x)} \\ \pause
                       &= (f(x_0) + f'(x_0)(x - x_0)) \cdot (\dfrac{1}{g(x_0)} + \dfrac{-g'(x_0)}{g(x_0)^2} (x - x_0)) \\ \pause
                       &= \dfrac{f(x_0)}{g(x_0)} + \dfrac{f'(x_0)}{g(x_0)}(x-x_0) + \dfrac{-f(x_0)g'(x_0)}{g(x_0)^2}(x - x_0) + O((x - x_0)^2) \\ \pause
                       &= \dfrac{f(x_0)}{g(x_0)} + \dfrac{f'(x_0) g(x_0) - f(x_0) g'(x_0)}{g(x_0)^2}(x - x_0)
  \end{align}
\end{frame}

\begin{frame}
  \frametitle{Quotient Rule V}

  Moral of the story: it's HARD to actually get the derivative of a quotient \pause so in no universe would the formula be as easy as $\dfrac{f'(x)}{g'(x)}$. \pause It is actually: \pause

  \begin{alertblock}{Quotient Rule}\index{derivative!quotient rule}
    \begin{equation}
      \dfrac{d}{dx} \left(\dfrac{f(x)}{g(x)}\right) = \dfrac{f'(x) g(x) - f(x)g'(x)}{g(x)^2}
    \end{equation}
  \end{alertblock}
\end{frame}

\begin{frame}
  \frametitle{Quotient Rule Example I}
  E.g. computing:
  \begin{equation}
    \dfrac{d}{dx} \dfrac{2xe^x + 3x^2}{x^3 + 2x + 1}
  \end{equation} \pause

  \begin{align}
    \dfrac{d}{dx} \dfrac{2xe^x + 3x^2}{x^3 + 2x + 1} &= \dfrac{(\dfrac{d}{dx} (2xe^x + 3x^2)) \cdot (x^3 + 2x + 1) - (2xe^x + 3x^2) \cdot (\dfrac{d}{dx}(x^3 + 2x + 1))}{(x^3 + 2x + 1)^2} \\ \pause
                                                     &= \dfrac{(\dfrac{d}{dx}(2xe^x) + \dfrac{d}{dx}(3x^2))\cdot (x^3 + 2x + 1) - (2xe^x + 3x^2) \cdot (3x^2 + 2)}{(x^3 + 2x + 1)^2} \\ \pause
                                                     &= \dfrac{(2e^x + 2xe^x + 6x)\cdot(x^3 + 2x + 1) - (2xe^x + 3x^2)\cdot (3x^2 + 2)}{(x^3 + 2x + 1)^2}
  \end{align}
\end{frame}

\section{Trigonometric Derivatives}

\begin{frame}
  \frametitle{Derivative of $\sin(x)$ I}

  Recall from our discussion on the Squeeze Theorem:
  \begin{gather}
    \lim\limits_{x \rightarrow 0} \dfrac{\sin(x)}{x} = 1 \\
    \lim\limits_{x \rightarrow 0} \dfrac{\cos(x) - 1}{x} = 0
  \end{gather} \pause

  What is this limit in terms of derivatives? \pause

  \begin{gather}
    \dfrac{d}{dx} \mid_{x = 0} \sin(x) = 1,
    \dfrac{d}{dx} \mid_{x = 0} \cos(x) = 0
  \end{gather}
\end{frame}

\begin{frame}
  \frametitle{Derivative of $\sin(x)$ II}

  What about the derivative function? \pause

  \begin{align}
    \dfrac{d}{dx} \sin(x) &= \lim\limits_{h \rightarrow 0} \dfrac{\sin(x + h) - sin(x)}{h} \\ \pause
                          &= \lim\limits_{h \rightarrow 0} \dfrac{\sin(x)\cos(h) + \sin(h)\cos(x) - \sin(x)}{h} \\ \pause
                          &= \lim\limits_{h \rightarrow 0} \dfrac{\sin(x)(\cos(h) - 1) + \cos(x) \sin(h)}{h} \\ \pause
                          &= \lim\limits_{h \rightarrow 0} \left(\sin(x) \dfrac{\cos(h) - 1}{h} + \cos(x) \dfrac{\sin(h)}{h}\right) \\ \pause
                          &= \sin(x) \lim\limits_{h \rightarrow 0} \left( \dfrac{\cos(h) - 1}{h} \right) + \cos(x) \lim\limits_{h \rightarrow 0} \left(\dfrac{\sin(h)}{h}\right) \\ \pause
                          &= \sin(x) \cdot 0 + \cos(x) \cdot 1 \pause = \cos(x)
  \end{align}
\end{frame}

\begin{frame}
  \frametitle{Derivative of $\sin(x)$ III}
  \begin{alertblock}{Derivative of $\sin(x)$}\index{derivative! of sin}
    \begin{equation}
      \dfrac{d}{dx} \sin(x) = \cos(x)
    \end{equation}
  \end{alertblock} \pause

  What about for $\cos(x)$?
\end{frame}

\begin{frame}
  \frametitle{Derivative of $\cos(x)$ I}

  \begin{align}
    \dfrac{d}{dx} \cos(x) &= \lim\limits_{h \rightarrow 0} \dfrac{\cos(x + h) - \cos(x)}{h} \\ \pause
                          &= \lim\limits_{h \rightarrow 0} \dfrac{\cos(x)\cos(h) - \sin(x)\sin(h) - \cos(x)}{h} \\ \pause
                          &= \cos(x) \lim\limits_{h \rightarrow 0} \left( \dfrac{\cos(h) - 1}{h} \right) - \sin(x) \lim\limits_{h \rightarrow 0} \left(\dfrac{\sin(h)}{h} \right) \pause = - \sin(x)
  \end{align} \pause

  \begin{alertblock}{Derivative of $\cos(x)$}\index{derivative! of cos}
    \begin{equation}
      \dfrac{d}{dx} \cos(x) = -\sin(x)
    \end{equation}
  \end{alertblock}
\end{frame}

\begin{frame}
  \frametitle{Derivative of $\tan(x)$}
  What about for $\tan(x)$? Do we need to do even more limit yuckiness? \pause

  No! Since $\tan(x) = \dfrac{\sin(x)}{\cos(x)}$ we can just use quotient rule! \pause

  \begin{align}
    \dfrac{d}{dx} \tan(x) &= \dfrac{d}{dx} \dfrac{\sin(x)}{\cos(x)} \\ \pause
                          &= \dfrac{\left(\dfrac{d}{dx}\sin(x)\right)\cdot \cos(x) - \sin(x) \cdot \left(\dfrac{d}{dx} \cos(x)\right)}{\cos(x)^2} \\ \pause
                          &= \dfrac{\cos(x) \cdot \cos(x) - \sin(x) \cdot (- \sin(x))}{\cos(x)^2} \\ \pause
                          &= \dfrac{\cos(x)^2 + \sin(x)^2}{\cos(x)^2} \\ \pause
                          &= \dfrac{1}{\cos(x)^2} \pause = \sec(x)^2
  \end{align}
\end{frame}

\begin{frame}
  \frametitle{Derivative of $\sin(2x)$}

  While we'll get an easier method to do this next lecture, we can also use trig rules to compute $\dfrac{d}{dx} \sin(2x)$: \pause

  \begin{align}
    \dfrac{d}{dx} \sin(2x) &= \dfrac{d}{dx}(2 \sin(x) \cos(x)) \\ \pause
                           &= 2 \dfrac{d}{dx} (\sin(x) \cos(x)) \\ \pause
                           &= 2 \cdot \left( \left(\dfrac{d}{dx} \sin(x)\right) \cdot \cos(x) + \sin(x) \cdot \left( \dfrac{d}{dx} \cos(x)\right) \right) \\ \pause
                           &= 2 \cdot (\cos(x) \cdot \cos(x) + \sin(x) \cdot (- \sin(x))) \\ \pause
                           &= 2 \cdot (\cos(x)^2 - \sin(x)^2) \\ \pause
                           &= 2 \cdot \cos(2x)
  \end{align}
\end{frame}

\begin{frame}
  \frametitle{Worksheet Time!}

  Let's get into groups (or if you really like working alone, please stay so I can be here to help :) ) and work on the worksheet!
\end{frame}

\begin{frame}
  \frametitle{End of Class!}

  That's all folks!
\end{frame}
\end{document}
